% Part: model-theory
% Chapter: models-of-arithmetic
% Section: standard-models

\documentclass[../../../include/open-logic-section]{subfiles}

\begin{document}

\olfileid[pt]{mod}{mar}{stm}
\olsection{Modelos padrão da aritmética}

A linguagem da aritmética~$\Lang{L_A}$ destina-se, evidentemente, a
falar de números, em particular dos números naturais. Por isso, ``o''
modelo padrão~$\Struct{N}$ é especial: é o modelo sobre o qual queremos
falar. Em lógica, porém, muitas vezes nos interessam apenas as
propriedades estruturais, e quaisquer duas !!{structure}s isomorfas
compartilham essas propriedades. Assim, podemos ser um pouco mais
liberais e chamar de ``padrão'' qualquer !!{structure} isomorfa
a~$\Struct{N}$.

\begin{defn}
Uma !!{structure} para $\Lang{L_A}$ é \emph{padrão} se for isomorfa
a~$\Struct{N}$.
\end{defn}

\begin{prop}
\ollabel{prop:standard-domain} Se !!a{structure}~$\Struct{M}$ é padrão,
então seu domínio é o conjunto dos valores dos numerais padrão, isto é,
\[
\Domain{M} = \Setabs{\Value{\num{n}}{M}}{n \in \Nat}
\]
\end{prop}

\begin{proof}
É claro que todo $\Value{\num{n}}{M} \in \Domain{M}$. Basta mostrar
que todo $x \in \Domain{M}$ é igual a $\Value{\num{n}}{M}$ para
algum~$n$. Como $\Struct{M}$ é padrão, ela é isomorfa
a~$\Struct{N}$. Suponha que $g\colon \Nat \to \Domain{M}$ seja um
isomorfismo. Então $g(n) = g(\Value{\num{n}}{N}) =
\Value{\num{n}}{M}$. Mas, para todo $x \in \Domain{M}$, existe~$n
\in \Nat$ tal que $g(n) = x$, pois $g$ é !!{surjective}.
\end{proof}

\begin{explain}
Se uma estrutura~$\Struct{M}$ para $\Lang{L_A}$ é padrão, todos os
elementos de seu !!{domain} podem ser nomeados pelos numerais padrão
$\num{0}$, $\num{1}$, $\num{2}$, \dots, isto é, pelos termos $\Obj{0}$,
$\Obj{0}'$, $\Obj{0}''$ etc. Isso, naturalmente, não significa que os
!!{element}s de $\Domain{M}$ \emph{sejam} os números, mas apenas que
podemos distingui-los da mesma maneira que distinguimos os números
em~$\Domain{N}$.
\end{explain}

\begin{prob}
Mostre que a recíproca de \olref[mod][mar][stm]{prop:standard-domain}
é falsa, isto é, dê um exemplo de !!a{structure}~$\Struct{M}$ com
$\Domain{M} = \Setabs{\Value{\num{n}}{M}}{n \in \Nat}$ que não seja
isomorfa a~$\Struct{N}$.
\end{prob}

\begin{prop}
\ollabel{prop:thq-standard}
Se $\Sat{M}{\Th{Q}}$ e $\Domain{M} = \Setabs{\Value{\num{n}}{M}}{n
  \in \Nat}$, então $\Struct{M}$ é padrão.
\end{prop}

\begin{proof}
Devemos mostrar que $\Struct{M}$ é isomorfa
a~$\Struct{N}$. Considere a função $g\colon \Nat \to \Domain{M}$
definida por $g(n) = \Value{\num{n}}{M}$. Pela hipótese, $g$ é
!!{surjective}. Ela também é !!{injective}: $\Th{Q} \Proves
\eq/[\num{n}][\num{m}]$ sempre que $n \neq m$. Portanto, como
$\Sat{M}{\Th{Q}}$, temos $\Sat{M}{\eq/[\num{n}][\num{m}]}$ sempre que
$n \neq m$. Logo, se $n \neq m$, então $\Value{\num{n}}{M} \neq
\Value{\num{m}}{M}$, isto é, $g(n) \neq g(m)$.

Também precisamos verificar que $g$ é um isomorfismo.
\begin{enumerate}
\item Temos $g(\Assign{\Obj{0}}{N}) = g(0)$, pois
  $\Assign{\Obj{0}}{N} = 0$. Pela definição de~$g$, $g(0) =
  \Value{\num{0}}{M}$. Mas $\num{0}$ é simplesmente $\Obj{0}$, e o
  valor de um termo que é !!a{constant} é dado pelo que a !!{structure}
  atribui a essa !!{constant}, isto é,
  $\Value{\Obj{0}}{M} = \Assign{\Obj{0}}{M}$. Assim, obtemos
  $g(\Assign{\Obj{0}}{N}) = \Assign{\Obj{0}}{M}$, como exigido.
\item $g(\Assign{\prime}{N}(n)) = g(n+1)$, pois $\prime$ em
  $\Struct{N}$ é a função sucessor em~$\Nat$. Então $g(n+1) =
  \Value{\num{n+1}}{M}$ pela definição de~$g$. Mas $\num{n+1}$ é o
  mesmo termo que $\num{n}'$, de modo que $\Value{\num{n+1}}{M} =
  \Value{\num{n}'}{M}$. Pela definição da função valor, isso é
  $= \Assign{\prime}{M}(\Value{\num{n}}{M})$. Como
  $\Value{\num{n}}{M} = g(n)$, obtemos $g(\Assign{\prime}{N}(n)) =
  \Assign{\prime}{M}(g(n))$.
\item $g(\Assign{+}{N}(n,m)) = g(n+m)$, pois $+$ em $\Struct{N}$ é a
  função adição em~$\Nat$. Então $g(n+m) = \Value{\num{n+m}}{M}$ pela
  definição de~$g$. Mas $\Th{Q} \Proves \num{n+m} = (\num{n} +
  \num{m})$, logo $\Value{\num{n+m}}{M} =
  \Value{\num{n}+\num{m}}{M}$. Pela definição da função valor, isso é
  $= \Assign{+}{M}(\Value{\num{n}}{M},\Value{\num{m}}{M})$. Como
  $\Value{\num{n}}{M} = g(n)$ e $\Value{\num{m}}{M} = g(m)$, obtemos
  $g(\Assign{+}{N}(n, m)) = \Assign{+}{M}(g(n), g(m))$.
\item $g(\Assign{\times}{N}(n, m)) = \Assign{\times}{M}(g(n), g(m))$:
  Exercício.
\item $\tuple{n,m} \in \Assign{<}{N}$ se, e somente se, $n < m$. Se
  $n < m$, então $\Th{Q} \Proves \num{n} < \num{m}$ e também
  $\Sat{M}{\num{n} < \num{m}}$. Portanto,
  $\tuple{\Value{\num{n}}{M}, \Value{\num{m}}{M}} \in
  \Assign{<}{M}$, isto é, $\tuple{g(n), g(m)} \in \Assign{<}{M}$. Se
  $n \not< m$, então $\Th{Q} \Proves \lnot \num{n} < \num{m}$ e,
  consequentemente, $\Sat/{M}{\num{n} < \num{m}}$. Logo, como antes,
  $\tuple{g(n), g(m)} \notin \Assign{<}{M}$. Em conjunto, obtemos:
  $\tuple{n,m} \in \Assign{<}{N}$ se, e somente se,
  $\tuple{g(n), g(m)} \in \Assign{<}{M}$.
\end{enumerate}
\end{proof}

\begin{explain}
A função~$g$ é a maneira mais óbvia de definir uma aplicação de
$\Nat$ no domínio de qualquer outra !!{structure}~$\Struct{M}$ para
$\Lang{L_A}$, pois toda $\Struct{M}$ desse tipo contém !!{element}s
nomeados por $\num{0}$, $\num{1}$, $\num{2}$ etc. Portanto, não
surpreende que, se $\Struct{M}$ torna verdadeiras ao menos algumas
afirmações básicas sobre os $\num{n}$ do mesmo modo que~$\Struct{N}$,
e se $g$ também é bijetiva, então $g$ seja um isomorfismo. De fato, se
$\Domain{M}$ não contém !!{element}s além dos nomeados pelos
$\num{n}$, esse é o único isomorfismo.
\end{explain}

\begin{prop}
\ollabel{prop:thq-unique-iso} Se $\Struct{M}$ é padrão, então a função
$g$ da demonstração de \olref{prop:thq-standard} é o único isomorfismo
de $\Struct{N}$ para~$\Struct{M}$.
\end{prop}

\begin{proof}
Suponha que $h\colon \Nat \to \Domain{M}$ seja um isomorfismo entre
$\Struct{N}$ e~$\Struct{M}$. Mostramos que $g = h$ por indução
em~$n$. Se $n = 0$, então $g(0) = \Assign{\Obj{0}}{M}$ pela definição
de~$g$. Mas, como $h$ é um isomorfismo, $h(0) =
h(\Assign{\Obj{0}}{N}) =\Assign{\Obj{0}}{M}$, portanto $g(0) = h(0)$.

Considere agora o caso $n+1$. Temos
\begin{align*}
  g(n+1) & = \Value{\num{n+1}}{M} \text{ pela definição de~$g$}\\
  & = \Value{\num{n}'}{M} \text{ pois $\num{n+1}\ident \num{n}'$}\\
  & = \Assign{\prime}{M}(\Value{\num{n}}{M}) 
    \text{ pela definição de $\Value{t'}{M}$}\\
  & = \Assign{\prime}{M}(g(n))  \text{ pela definição de~$g$}\\
  & = \Assign{\prime}{M}(h(n)) \text{ pela hipótese de indução}\\
  & = h(\Assign{\prime}{N}(n)) \text{ pois $h$ é um isomorfismo}\\
  & = h(n+1)
\end{align*}
\end{proof}

\begin{explain}
Para qualquer conjunto !!{denumerable}~$M$, existe !!a{bijection}
entre $\Nat$ e $M$, de modo que todo conjunto~$M$ desse tipo pode ser
o !!{domain} de um modelo padrão~$\Struct{M}$. Na verdade, depois de
escolher um objeto $z \in M$ e uma função adequada $s$ como
$\Assign{\Obj{0}}{M}$ e $\Assign{\prime}{M}$, as interpretações de
$+$, $\times$ e $<$ já ficam determinadas. Apenas funções
$s\colon M \to M \setminus \{z\}$ que sejam ao mesmo tempo
!!{injective} e !!{surjective} são adequadas como
$\Assign{\prime}{M}$ em um modelo padrão. A imagem de $s$ não pode
conter~$z$, pois, caso contrário, $\lforall[x][\eq/[\Obj 0][x']]$
seria falsa. Essa !!{sentence} é verdadeira em~$\Struct{N}$, portanto
$\Struct{M}$ também deve torná-la verdadeira. A função~$s$ deve ser
!!{injective}, pois a função sucessor~$\Assign{\prime}{N}$ em
$\Struct{N}$ o é, e o fato de $\Assign{\prime}{N}$ ser !!{injective}
é expresso por !!a{sentence} verdadeira em~$\Struct{N}$. Ela deve ser
!!{surjective}, pois, caso contrário, existiria algum $x \in M
\setminus \{z\}$ que não estaria na imagem de~$s$, isto é, a
!!{sentence} $\lforall[x][(\eq[x][\Obj 0] \lor
\lexists[y][\eq[y'][x]])]$ seria falsa em~$\Struct{M}$---mas é
verdadeira em~$\Struct{N}$.
\end{explain}


\end{document}
